\chapter{Cronograma e Formação Acadêmica}
\label{chap:cronograma}

O presente capítulo apresenta o planejamento temporal das atividades do doutorado, bem como a organização das disciplinas obrigatórias e o aproveitamento de créditos já cursados em outros programas.

\section{Cronograma de Execução}

Esta seção detalha o cronograma a ser seguido no Programa de Pós-Graduação em Modelagem Matemática e Computacional sob orientação do Prof. Dr. Alisson Marques da Silva no cumprimento das exigências necessárias à obtenção do título de Doutor em Modelagem Matemática e Computacional. Para tanto, na Tabela~\ref{tab:cronograma} apresenta-se o cronograma de atividades proposto para realização deste projeto, considerando o prazo máximo de 8 semestres para sua integralização.

A pesquisa será desenvolvida ao longo de 48 meses (8 semestres), dividida nas seguintes etapas:

\begin{table}[h]
\centering
\caption{Cronograma de Atividades da Pesquisa}
\label{tab:cronograma}
\begin{tabular}{|l|c|c|c|c|c|c|c|c|}
\hline
\textbf{Atividade / Semestre} & \textbf{1} & \textbf{2} & \textbf{3} & \textbf{4} & \textbf{5} & \textbf{6} & \textbf{7} & \textbf{8} \\
\hline
Integralização dos créditos & X & & & & & & & \\
\hline
Revisão Bibliográfica & X & X & X & X & X & X & & \\
\hline
Proposição de algoritmos & & X & X & X & X & & & \\
\hline
Implementação e testes & & & X & X & X & X & & \\
\hline
Validação experimental & & & & X & X & X & X & \\
\hline
Qualificação & & & & X & & & & \\
\hline
Publicações científicas & & & X & X & X & X & X & \\
\hline
Escrita e defesa da tese & & & & & & X & X & X \\
\hline
\end{tabular}
\end{table}

\subsection{Descrição das Etapas}

\begin{itemize}
\item \textbf{Semestre 1:} Integralização de créditos e início da revisão bibliográfica
\item \textbf{Semestres 2-3:} Aprofundamento da fundamentação teórica e início da proposição de algoritmos
\item \textbf{Semestres 4-5:} Desenvolvimento e implementação dos algoritmos neuro-fuzzy evolutivos
\item \textbf{Semestre 4:} Exame de qualificação
\item \textbf{Semestres 6-7:} Validação experimental e análise comparativa
\item \textbf{Semestres 6-8:} Redação da tese e preparação para defesa
\item \textbf{Semestres 3-7:} Produção contínua de publicações científicas
\end{itemize}

\section{Disciplinas do Doutorado e Planejamento}

Esta seção detalha as disciplinas obrigatórias do doutorado, sua situação atual e como elas se encaixam no planejamento acadêmico do curso.

\subsection{Disciplinas do Doutorado (PPGMMC/CEFET-MG)}

O doutorado em Modelagem Matemática e Computacional (PPGMMC/CEFET-MG) exige o cumprimento de disciplinas obrigatórias e optativas, além de atividades de pesquisa e defesa. Abaixo, estão listadas as disciplinas já cursadas (inclusive como aluno especial) e aquelas ainda pendentes:

\begin{table}[h]
\centering
\caption{Disciplinas do Doutorado – Cursadas e Pendentes}
\label{tab:disciplinas_doutorado}
\resizebox{\textwidth}{!}{%
\begin{tabular}{|l|l|l|c|c|c|}
\hline
\textbf{Disciplina} & \textbf{Professor} & \textbf{Situação} \\
\hline
Algoritmos e Estruturas de Dados & Thiago de Souza Rodrigues & Obrigatória/Cursando (Isolada) \\
\hline
Inteligência Computacional & Alisson Marques da Silva & Obrigatória/Cursando (Isolada) \\
\hline
Álgebra Linear & -- & Obrigatória/A cursar \\
\hline
Princípios de Modelagem Matemática & -- & Obrigatória/A cursar \\
\hline
Engenharia de Software & -- & Obrigatória/A cursar \\
\hline
Planejamento e Análise Estatística de Experimentos & -- & Obrigatória/A cursar \\
\hline
Desenvolvimento de Projeto de Tese I & -- & Obrigatória/A cursar \\
\hline
Desenvolvimento de Projeto de Tese II & -- & Obrigatória/A cursar \\
\hline
Desenvolvimento de Projeto de Tese III & -- & Obrigatória/A cursar \\
\hline
Elaboração de Projeto de Tese & -- & Obrigatória/A cursar \\
\hline
Defesa de Tese & -- & Obrigatória/A cursar \\
\hline
Exame de Qualificação - Doutorado & -- & Obrigatória/A cursar \\
\hline
\end{tabular}%
}
\end{table}


\section{Formação Acadêmica e Disciplinas Cursadas}

Esta seção apresenta a formação acadêmica complementar realizada em diferentes programas de pós-graduação, demonstrando o embasamento teórico multidisciplinar para o desenvolvimento desta pesquisa.\newline
\textbf{Mestrado Profissional em Modelagem Computacional e Sistemas} pela Universidade Estadual de Montes Claros (UNIMONTES) -- PPGMCS.\newline
Na UFMG, o curso de doutorado nessa área se chama \textbf{Programa de Pós-graduação em Ciência da Computação} (PPGCC), oferecido pelo Departamento de Ciência da Computação (DCC), vinculado ao Instituto de Ciências Exatas (ICEx).


\subsection{Disciplinas Cursadas no Mestrado Profissional (UNIMONTES)}

No Mestrado Profissional em Modelagem Computacional e Sistemas da Universidade Estadual de Montes Claros (UNIMONTES) -- PPGMCS, foram cursadas disciplinas fundamentais que forneceram a base teórica para o desenvolvimento de sistemas inteligentes, modelagem matemática e métodos de detecção de anomalias:

\begin{table}[ht]
\centering
\caption{Disciplinas Cursadas no Programa de Mestrado Profissional (UNIMONTES)}
\label{tab:disciplinas_mestrado}
\resizebox{\textwidth}{!}{%
\begin{tabular}{|l|l|l|c|}
\hline
\textbf{Semestre/Ano} & \textbf{Disciplina} & \textbf{Professor} & \textbf{Nota} \\
\hline
2016/2 & Métodos Matemáticos & Prof. Dr. Gustavo Foscolo de Moura Gomes & 98 \\
\hline
2016/2 & Detecção e Diagnóstico de Falhas em Sistemas Dinâmicos & Prof. Dr. Marcos Flávio Silveira Vasconcelos D'Angelo & 91 \\
\hline
2016/2 & Processos Estocásticos & Prof. Dr. Nilson Luiz Castelucio Brito & 90 \\
\hline
2017/1 & Projeto e Análise de Algoritmos & Prof. Dr. René Rodrigues Veloso & 80 \\
\hline
2017/1 & Inteligência Computacional & Prof. Dr. Marcos Flávio Silveira Vasconcelos D'Angelo & 88 \\
\hline
2017/2 & Visão Computacional & Prof. Dr. Antonio Vilson Vieira & 86 \\
\hline
\end{tabular}%
}
\end{table}

\subsection{Disciplinas Isoladas (UFMG)}

Para complementar a formação e aprofundar conhecimentos específicos relacionados ao projeto, foram cursadas disciplinas isoladas na Universidade Federal de Minas Gerais:

\begin{table}[h]
\centering
\caption{Disciplinas Isoladas Cursadas na UFMG}
\label{tab:disciplinas_ufmg}
\resizebox{\textwidth}{!}{%
\begin{tabular}{|l|l|l|c|}
\hline
\textbf{Semestre/Ano} & \textbf{Disciplina} & \textbf{Professor} & \textbf{Nota} \\
\hline
2021/2 & Mineração de Dados & Prof. Dr. Wagner Meira Júnior & 86,0 \\
\hline
2021/2 & Tópicos Especiais em Ciência da Computação & Prof. Dr. Cristiano Arbex Valle & 95,0 \\
& \textit{(Finanças Quantitativas e Gerenciamento de Risco)} & & \\
\hline
2021/1 & Tópicos Especiais em Ciência da Computação & Prof. Dr. Erickson Rangel do Nascimento & 75,0 \\
& \textit{(Visão Computacional)} & & \\
\hline
2021/1 & Tópicos Especiais em Ciência da Computação & Profa. Dra. Raquel Cardoso de Melo Minardi & 99,0 \\
& \textit{(Visualização de Dados)} & & \\
\hline
\end{tabular}%
}
\end{table}


Esta formação abrangente, aliada ao cronograma estruturado de pesquisa, estabelece as condições ideais para o desenvolvimento bem-sucedido da investigação em detecção de fraudes utilizando sistemas neuro-fuzzy evolutivos aplicados ao contexto proposto. O aproveitamento de disciplinas já cursadas permitiu reduzir o tempo dedicado à formação básica, possibilitando maior foco no desenvolvimento, implementação e validação dos algoritmos propostos. Além disso, a experiência adquirida em áreas como modelagem matemática, inteligência computacional, análise de dados e mineração de informações é fundamental para a condução de uma pesquisa inovadora, robusta e alinhada às demandas atuais do setor. Dessa forma, o planejamento acadêmico e científico apresentado neste capítulo garante a viabilidade e a relevância do projeto, contribuindo para avanços significativos na detecção de fraudes e para a formação qualificado.
